\chapter{Supported anchored defects, exact cycle periods, and an
alphabet-exit certificate}
{\small\sffamily Source: \nolinkurl{repository/collatz_reconstruction/research_program/anchored_defect_20260914/note.md}.\par}


\textbf{14 September 2026. Collatz workbench continuation.}

Source: \nolinkurl{KokunoYumeto/collatz-workbench}, commit
\nolinkurl{6a30f1dc23110fae0abbe3b3bd5e60072a842385}, PR \#2. All
inherited files remain unchanged. The proofs below use the actual
original odd-return map and the preserved integral cycle and
residual-source constructions. The finite certificate is executed by
\nolinkurl{verify.py}; source preservation and all ten
ordinary/optimized executions are checked by \nolinkurl{replay.py}. No
independent external mathematical review, Lean verification, or global
convergence result is asserted.

\section{1. What the Split-Zero technique contributes, with an exact
comparison}

For a ring R, use the supplied support extension G(R)=R disjoint union
\{tau\}. The external tau is the additive identity and multiplicative
absorber. Operations on supported elements are their original ring
operations. Thus e=0\_R is a supported zero, with e distinct from tau.
There are explicit semiring maps

\SourceMath{\pi:G(R)\longrightarrow R,\qquad
 \chi:G(R)\longrightarrow\mathbb B,}{}

where pi fixes R and sends tau to 0, while chi sends tau to 0 and every
element of R, including e, to 1. The Boolean operations are OR and AND.
Direct case checking shows that both maps preserve addition and
multiplication. The map

\SourceMath{(\pi,\chi):G(R)\longrightarrow R\times\mathbb B}{S1}

is an isomorphism onto the subsemiring

\SourceMath{\{(0,0)\}\ \cup\ \{(r,1):r\in R\}.}{}

Its inverse sends (0,0) to tau and (r,1) to the supported r. Both
compositions are identities. The zero fiber of pi alone is \{tau,e\};
every nonzero fiber is a singleton. This also specifies exactly how
ordinary coefficient calculations with a retained Boolean support
coordinate implement the supplied technique. No exclusivity or priority
claim for these elementary operations is made.

Here R=Z or Q. The exponent word, source lattice, forcing vector and
graph edges are additional retained data, not information reconstructed
from chi alone. \nolinkurl{Supported} implements (S1) over Q and tests
both zero values, cancellation, multiplication and distributivity.

\subsection{A zero-forcing row is still an actual relation}

For the one-letter word (2), the anchored cycle equation is (4-3)y=0.
Its complex is {[}Z --1--\textgreater{} Z{]}, with both cohomology
groups zero. Deleting the zero-forcing row produces {[}Z
--\textgreater{} 0{]}, whose H\^{}0 is Z. The exact comparison is
q\^{}0=id and q\^{}1=0. Its kernel complex is {[}0 --\textgreater{}
Z{]}. In the associated long exact sequence the connecting map from
H\^{}0({[}Z--\textgreater0{]}) to H\^{}1({[}0--\textgreater Z{]}) is
multiplication by 1, computed by lifting the source generator and
applying the original differential. The extra kernel created by row
deletion is therefore explicitly accounted for. The supported zero never
licensed deletion of that differential.

\subsection{The coefficient map to the complex analytic setting}

The inclusion Z -\textgreater{} C induces G(Z) -\textgreater{} G(C),
taking tau to tau and the supported integer zero to the supported
complex zero. Through (S1), this is the restriction of inclusion times
id\_B, so it preserves the two zero labels injectively. For the original
return complex {[}Z --D--\textgreater{} Z{]}, coefficient extension
gives {[}C --D--\textgreater{} C{]}. At D nonzero, its contracting
homotopy is multiplication by 1/D in degree one. Thus the induced H\^{}1
map is the explicitly zero map Z/DZ -\textgreater{} 0, even though the
support coordinates remain present. For (1,3), D=7 and the forcing class
{[}5{]} is nonzero integrally; its complex primitive is 5/7. This proves
the precise loss under this coefficient change. We retain the original
integral complex rather than substitute its complex-coefficient
contraction.

The peer Zeta guide currently studies spectral cohomology and original
period/Gram calculations. That guide provides context, not an arithmetic
estimate used here. No identification of its theta complex with the
Collatz cycle complex is asserted. The coefficient maps just specified
are the comparisons used in this note.

\section{2. Original forcing and anchored primitive}

Let

\SourceMath{T(n)=\frac{3n+1}{2^{a(n)}},\quad a(n)=\nu_2(3n+1),\qquad
 X=\{1,3,5,\ldots\}.}{}

For an original nonempty exponent word p=(a\_1,\ldots,a\_m), retain

\SourceMath{A_j=\sum_{i=1}^j a_i,\quad A_0=0,\quad
 U=2^{A_m},\quad L=3^m,\quad D=U-L,}{}
\SourceMath{C=\sum_{i=1}^m3^{m-i}2^{A_{i-1}},\qquad
 F_p(n)=(Ln+C)/U.}{}

The preserved cycle differential and functional are

\SourceMath{(B_px)_i=2^{a_i}x_{i+1}-3x_i,\quad x_{m+1}=x_1,}{}
\SourceMath{f_i=3^{m-i}2^{A_{i-1}},\quad fB_p=D e_1^T.}{A1}

The predecessor proves H\^{}0=0 and H\^{}1=Z/DZ with {[}b{]} sent to
{[}fb{]}. The actual cyclic forcing is b=1, with marked coordinate
{[}C{]}. Its rational primitive is x\_i=C(p\_i)/D, for the retained
rotations p\_i. Integrality, positivity, and the actual valuation
equations remain required to obtain a positive integer orbit.

The affine coordinate map y=x-1 has inverse x=y+1. Substitution gives

\SourceMath{B_py=b_p,\qquad (b_p)_i=4-2^{a_i},\qquad fb_p=C-D.}{A2}

This is a bijection of the two specified solution sets, not a claim that
translation is a linear chain map. Its forcing difference is the
original boundary -B\_p1. Consequently

\SourceMath{[b_p]=[1],\qquad [C-D]=[C]\text{ in }\mathbb Z/D\mathbb Z.}{A3}

The same integral class is retained. What changes is its specified
primitive and the exact observable below. The trivial word repetitions
(2)\^{}m have y=0. Conversely y=0 forces every right-hand side
4-2\^{}\{a\_i\} to vanish, so every letter is 2 and x=1. An actual
nontrivial positive cycle never meets 1, since T(1)=1; therefore all its
x\_i are at least 3 and y\_i are positive even integers.

The block maps from
\nolinkurl{../residual_splice_20260914/integral_blocks.md} still apply.
In that notation the anchored right-hand side maps to Rb\_p, and the
inverse is y=S\_0z+Hb\_p, followed by x=y+1. Dropping Hb\_p does not
give this affine inverse. Equations (A2)-(A3) make the comparison with
the existing marked obstruction literal.

\section{3. An anchored period character with its complete zero domain}

For n\textgreater1 define

\SourceMath{\eta(n)=\frac{2^{a(n)}(T(n)-1)}{3(n-1)}
 =1+\frac{4-2^{a(n)}}{3(n-1)}.}{P1}

This is a nonnegative rational value. A step into 1 gives eta=0,
retained as a supported zero. At n=1 the ratio is not defined: its
retained equation is 4*0=3*0+0, with exponent 2 and zero forcing. No 0/0
is assigned a value.

For an actual finite path starting above 1 and stopped no later than its
first arrival at 1, multiplication of the exact edge equations gives

\SourceMath{\prod_{j=0}^{m-1}\eta(T^j(n))
 =\frac{2^{A_m}}{3^m}\frac{T^m(n)-1}{n-1}.}{P2}

The identity includes terminal value 1: both sides are then supported
zero. Its proof cancels only the positive intermediate factors. No
canceled intermediate factor is zero.

On the subgraph with both endpoints above 1, all eta values are
positive. Let C\_1 be the free abelian group on its actual edges, C\_0
the free abelian group on its vertices, and partial(e\_n)=v\_n-v\_T(n).
Define multiplicative homomorphisms

\SourceMath{\rho(v_n)=n-1,\qquad \lambda(e_n)=2^{a(n)}/3.}{}

They extend by products with the original integer chain coefficients.
Their exact relation is

\SourceMath{\eta(c)=\lambda(c)\rho(\partial c)^{-1}.}{P3}

Thus eta and lambda differ by the specified vertex coboundary, and
define the same class in multiplicative degree-one graph cohomology. In
particular every actual closed cycle c satisfies

\SourceMath{\eta(c)=\lambda(c)=2^A/3^m.}{P4}

The kernel of lambda consists exactly of finite chains with both sum
c\_n a(n)=0 and sum c\_n=0, by unique prime factorization. The kernels
of rho and eta are likewise the simultaneous integer prime-valuation
equations of their displayed rational products. These describe the
fibers as cosets of the stated kernels; no unspecified spectral quotient
is used.

The predecessor's retractions preserve closed chains literally,
F(c)=c.~Formula (P4) is therefore unchanged on every actual nontrivial
cycle under those retractions. We do not extend the positive-group
character to a general retraction path that enters 1: that path has the
supported zero endpoint in (P2), not an invertible character value.

\subsection{Theorem P. A cycle inequality controlled by the original
exponent-1 positions}

For a nontrivial positive integer cycle with every vertex at least
s\textgreater=3, and k occurrences of exponent 1,

\SourceMath{\boxed{\quad
 1<\frac{2^A}{3^m}
 \le \left(\frac{3s-1}{3s-3}\right)^k.
 \quad}}{P5}

\textbf{Proof.} The ordinary cycle equations give
\nolinkurl{2^A=product(3+1/x_i)>3^m.} By (P1), an exponent-1 edge has
factor 1+2/(3(x\_i-1)), at most (3s-1)/(3s-3). An exponent-2 edge has
factor exactly 1: its zero anchored forcing remains a supported relation
and its original exponent still contributes to A. Every exponent at
least 3 has a factor strictly between 0 and 1, since the next cycle
vertex exceeds 1. Bound those latter factors above by 1, retain the
original k factors, and use (P4). This proves (P5). In particular k
cannot be zero. No independence, asymptotic estimate, or probabilistic
coefficient model occurs. QED.

For a fixed contracting pair (m,A), (P5) is an exact integer test

\SourceMath{2^A(3s-3)^k\le3^m(3s-1)^k.}{P6}

Its right-to-left ratio strictly decreases with integer s\textgreater=2.
An accepted integer and rejection of the next one therefore certify its
exact ceiling. For (m,A,k)=(971,1539,403), the ceiling is 274546. The
prior unanchored ceiling was 330749. Both endpoint comparisons are
checked on integers.

\section{4. The complete next exponent partition and its exact parameter
maps}

Keep a legal prefix q, including the empty prefix, with the original
chart

\SourceMath{n=r+2Ut,\quad x=F_q(n)=u+2Lt,\quad Uu=Lr+C.}{C1}

For the empty word use L=U=1, C=0, r=u=1. Every prefix source under
consideration is a positive odd integer. The next three cases are

\SourceMath{a=1,\qquad a=2,\qquad a\ge3.}{C2}

The first two cases are precisely the original legal cylinders of q1 and
q2. The last case is exactly

\SourceMath{t=t_0+4v,\qquad
 t_0\equiv-\frac{3u+1}{2}(3L)^{-1}\pmod4,\quad0\le t_0<4,}{}
\SourceMath{n=r+2Ut_0+8Uv.}{C3}

The full exponent in this last case remains recoverable:

\SourceMath{a=\nu_2(3u+1+6Lt).}{C4}

Indeed, 3x+1 is positive and even. Its valuation is at least 3 precisely
when (3u+1)/2+3Lt is divisible by 4. Since 3L is odd, (C3) is its
complete solution. Its inverse is v=(n-r-2Ut\_0)/(8U), and t=(n-r)/(2U).
This proves both the partition and its full label recovery. The lower
bound a\textgreater=3 is not a substitution a=3.

For each of the two exact-exponent children, the original descent test
remains

\SourceMath{F_p(n)<n\quad\Longleftrightarrow\quad D_p>0\text{ and }D_p n>C_p.}{C5}

At a node that has not previously descended, (C5) splits the current
source interval into a descent part above floor(C\_p/D\_p) and a
retained part at or below it. For D\_p\textless=0 there is no descent
part. Empty parts keep their labels and original equations.

\subsection{A synchronized congruence has its own retained map}

The cover also imposes n=c mod M. To intersect this with n=r mod Q, put
g=gcd(Q,M). Compatibility is g divides c-r. In the compatible case
choose

\SourceMath{t_0\equiv ((c-r)/g)(Q/g)^{-1}\pmod{M/g},\quad 0\le t_0<M/g,}{}
\SourceMath{\alpha=r+Qt_0,\quad\Delta=Q(M/g),\qquad n=\alpha+\Delta z.}{C6}

Both original residue equations are recovered by this substitution.
Conversely every joint solution has this form, because its original
t=(n-r)/Q solves the displayed linear congruence. The inverses are
z=(n-alpha)/Delta and t=(n-r)/Q. Intersecting with the retained interval
{[}l,h{]} gives exactly

\SourceMath{\left\lceil\frac{l-\alpha}{\Delta}\right\rceil
 \le z\le
 \left\lfloor\frac{h-\alpha}{\Delta}\right\rfloor.}{C7}

The code retains incompatible residue defects, compatible empty
intervals, the original r,Q,c,M, and both parameter maps. It does not
replace the source progression by a uniformly sampled interval.

\section{5. An actual alphabet-exit map, not a convergence certificate}

Define the graph G\_12 on the actual positive odd vertices by retaining
exactly the original edges e\_n for which a(n) is 1 or 2. Let E be the
actual vertex set with a(n)\textgreater=3. The inclusion of G\_12 into
the original orbit graph is identity on vertices and on the retained
edges and commutes with partial. Every E-vertex is present but has no
outgoing edge in G\_12.

An original positive integral solution for a word in \{1,2\} maps its
position circle into G\_12: position i goes to x\_i and its edge goes to
e\_x\_i. The original equations establish those edge memberships and the
boundary identity. The block-circle expansion of the predecessor sends
each block edge to the sum of its original edges, so its fundamental
cycle maps to the same full cycle chain. Repetitions keep their integer
multiplicities.

There is also the explicit relative chain map from the complex relative
to (v\_1,e\_1) to the one additionally relative to E: q\_1 is the
identity, and q\_0 kills the E-vertices and fixes the others. Its
degree-zero kernel is the free group on E; its degree-one kernel is
zero. Because both complexes have no degree-two chains, q induces an
injection on H\_1. In particular, a nontrivial \{1,2\} cycle stays a
nonzero cycle on its original edges.

A prefix in (C3) gives a path of allowed original edges from n to z in
E, with boundary v\_n-v\_z. This becomes a boundary certificate for v\_n
in the exit-relative complex. It remains v\_n-v\_z in the original
graph. There is no assertion that z reaches 1.

For a proposed cycle minimum n, either of two original facts excludes
membership: a path descends below n, or a path of allowed edges reaches
E. The former contradicts minimality; the latter contradicts the
proposed \{1,2\} word. This is the exact relationship used by the finite
certificate. It does not erase the possibility of a cycle having an
exponent at least 3.

\section{6. Discharge of the complete 403-rise sector}

The preserved residual-splice certificate proves an actual first descent
for every odd source from 3 through 99779. Strong induction proves that
all those sources reach 1. Its existing all-length theorem excludes at
most 402 exponent-1 positions. The new calculation keeps that common
source bound unchanged.

\subsection{Theorem Q. No nontrivial positive cycle has at most 403
exponent-1 positions in its primitive odd-return word}

\textbf{Proof.} Its minimum is at least s=99781 by the replayed
predecessor certificate. The original product and exponent identities
give

\SourceMath{(4s)^m\le2^k(3s+1)^m,\qquad
 A=2m-k+E,\quad E=\sum_{a_i\ge3}(a_i-2)\ge0.}{Q1}

For k\textless=403, exact integer comparison at m=972 and monotonicity
of (4s)/(3s+1)\textgreater1 force m\textless=971. For every length 1
through 970, the exact comparison is 2\^{}bitlength(3\^{}m) * s\^{}m
\textgreater{} (3s+1)\^{}m; thus the least power of 2 greater than
3\^{}m already exceeds the upper cycle bound. These 970 comparisons are
checked independently by a least-power-of-two loop. At m=971 the only
admissible exponent sum is A=1539; the next power fails the upper bound.
Then E=k-403\textgreater=0, giving k=403 and E=0. Thus the word consists
of 403 ones and 568 twos.

Theorem P bounds its minimum by 274546. Its outgoing exponent must be 1,
since an exponent 2 would descend at a source greater than 1. Its
incoming exponent must be 2, since exponent 1 strictly increases a
positive source. The first fact gives n=3 mod4. The second gives 4n=3x+1
for its original predecessor x, hence n=1 mod3. The CRT inverse (C6)
identifies these two conditions with n=7 mod12. The possible minima are
exactly

\SourceMath{n=99787,99799,\ldots,274543,}{Q2}

with 14564 elements. No word-order choice has been discarded: every
original cycle in the sector maps its minimum to a member of (Q2).

Starting from this precise source, the complete partitions (C2)-(C7)
give a finite proof tree. At every nonempty node all previous actual
returns have been at least the original source. The tree retains the
original word, congruences, all coefficients, and all empty child masks.
This instance uses the code's excess budget 0; the positive-budget
comparison is proved in the companion note. It closes after 2957
branching nodes; no nonempty prefix exceeds 32 allowed returns and no
OPEN leaf remains. Of its 8871 terminal labels, 7721 have empty integer
fibers and keep their present labels. Exactly 1150 labels are occupied:
1133 actual alphabet-exit tails and 17 first-descent cells. Their
cardinalities sum to

\SourceMath{13399\text{ exit sources}+1165\text{ descent sources}=14564.}{Q3}

The certificate checks disjointness and conservation at every node and
checks every terminal fiber by independently enumerating the ORIGINAL
dyadic progression and imposing n=7 mod12, without using the generator's
CRT inverse. A separate repeated-division interpreter reads every n in
(Q2), verifies its full word and endpoint classification, and
accumulates the original path boundary. Every one descends or exits the
allowed alphabet. Section 5 excludes both cases for the proposed cycle
minimum. This proves the theorem. QED.

This is a computer-assisted finite certificate combined with the
displayed all-length inequalities. It advances the self-contained
workbench bound, not a claimed published cycle record. It does not
depend on independent random exponents or on the distributional
estimates in the earlier notes.

The proof-tree generator accepts arbitrary finite depth caps. A cap
reached before closure creates an explicit OPEN leaf; the verifier
rejects its use as a proof of this exclusion. Other rejected controls
include an erased empty label, a forged original word, a deleted
occupied fiber, and an alphabet exit incorrectly promoted to descent.

\section{7. What remains and how to reproduce the work}

The companion \nolinkurl{turns_and_unit_excess.md} proves the ordered
adjacent-turn period identity, retains the exponent-3 negative forcing,
and completely discharges the 404-rise sector too. Its two source
classes have different exact incoming-edge constraints. The next
405-rise sector remains at m=971, A=1539 with total excess 2: either two
threes and 564 twos, or one four and 565 twos, together with the 405
ones. These original letters and their placement remain part of the
task; a zero-budget alphabet exit is not a certificate for that larger
fiber.

For the nonperiodic direction, (P2) retains the endpoint displacement as
well as both clocks. No uniform bound on that endpoint or on the
residual integer-supported source has been proved here. The old residual
roots and their degree-zero obstruction remain. The alphabet-exit
operation is confined to its specified word fiber and does not supply a
source-complete contraction of the original graph.

Run from the repository root:

\begin{Verbatim}[breaklines=true,breakanywhere=true,fontsize=\scriptsize]
python -B collatz_reconstruction/research_program/anchored_defect_20260914/replay.py
\end{Verbatim}

To materialize every tree node and every retained zero label:

\begin{Verbatim}[breaklines=true,breakanywhere=true,fontsize=\scriptsize]
python -B collatz_reconstruction/research_program/anchored_defect_20260914/anchored_cycles.py --full-cover --output cover.json
\end{Verbatim}

The full-tree hash, independent witness hash, example actual paths, all
period endpoints, and rejected controls are retained in
\nolinkurl{verification.json}. The generator, independently enumerated
fiber checks, and repeated-division interpreter are supplied in full.

\section{Sources and attribution}

\textbf{{[}S{]}} The supplied Split-Zero support extension, used in the
pinned Collatz \nolinkurl{split_zero_history_20260913/note.tex},
subsection ``Support, specialization, and what each forgets.'' Equation
(S1) proves the representation used here. User direction is retained as
Split-Zero cohomology and the related exact-map technique; an earlier
speech-to-text acronym is not used as a mathematical attribution.

\textbf{{[}P{]}} Pinned Collatz
\nolinkurl{cycle_relative_cohomology_20260914/note.md}, Sections 1, 5
and 6: original marked cycle complex, actual relative graph and chain
retractions. \nolinkurl{residual_splice_20260914/note.md}, Theorem 4 and
the executed 99779 source bound; \nolinkurl{integral_blocks.md},
equations B5-B13 and its integral affine inverse. All are retained
unchanged and replayed.

\textbf{{[}Z{]}} The connected peer Zeta
\nolinkurl{CURRENT_RESEARCH.md}, blob
\nolinkurl{9eefdcb14ca6fa22efb29fd60a01429512bfbadc}, read 14 September
2026, opens with original Gamma volumes, intrinsic sums, period Gram
matrices and section-corrected residuals. This current programme context
is not a premise of Theorems P or Q. No Zeta file is changed.

\textbf{{[}R{]}} O. Rozier and C. Terracol, \emph{Paradoxical behavior
in Collatz sequences}, arXiv:2502.00948v5, 17 May 2026, Introduction and
equation (2), \nolinkurl{https://arxiv.org/html/2502.00948v5}. The
shortened-clock affine expression maps to ours by expanding each
odd-return a into one odd shortened step followed by a-1 even steps:
shortened length is A, odd count is m, and the remainder is C/U. This
explicit comparison places the coefficient/offset calculation in the
existing literature. Their conjectural global assertions and numerical
bounds are not premises here. No priority claim for affine coding, cycle
products, or elementary cohomological algebra is made.
